\section{Problème Métier et Contexte}

\subsection{Enjeux et vulnérabilités des infrastructures métropolitaines}
La métropole moderne de Novapolis, à l'instar de tout grand centre urbain en forte croissance, fait face à des défis structurels majeurs en matière de gestion des risques et de planification opérationnelle. La complexité croissante de son tissu urbain – caractérisé par une imbrication dense de zones résidentielles, de pôles industriels, d'activités tertiaires et d'infrastructures vitales – expose la collectivité à des vulnérabilités systémiques. 

Parmi les entités les plus critiques figurent les infrastructures de santé (hôpitaux généraux et spécialisés), les services de sécurité civile (casernes de pompiers), les établissements d'enseignement accueillant du public et les centres névralgiques de télécommunication. En situation de crise (accident industriel majeur, catastrophe naturelle, congestion routière critique ou panne d'infrastructure numérique), la défaillance d'un seul de ces maillons peut paralyser l'ensemble de la chaîne de secours métropolitaine.

\subsection{Les trois axes fondamentaux de l'étude}
Pour répondre à ces exigences de sûreté et de continuité d'activité, les autorités locales de Novapolis ont mandaté une évaluation approfondie reposant sur la modélisation mathématique et la théorie des graphes. Ce projet intégrateur s'articule autour de trois questions fondamentales :

\begin{enumerate}
    \item \textbf{Accessibilité spatiale et équité des services d'urgence :} 
    Quels quartiers, îlots d'habitation ou établissements sensibles ne sont pas suffisamment connectés ou accessibles depuis les casernes de pompiers ou les centres hospitaliers ? Il s'agit d'identifier les zones d'ombre topologiques où les délais d'intervention dépassent les seuils acceptables, garantissant ainsi une équité territoriale face au risque.
    
    \item \textbf{Optimisation dynamique des itinéraires de secours :} 
    Quels chemins recommander aux véhicules d'intervention lorsque certaines voies routières stratégiques sont saturées en période de pointe, endommagées ou totalement obstruées par un sinistre ? Le modèle doit dépasser la simple notion de plus court chemin statique pour intégrer des scénarios de dégradation, évaluer les retards induits (absolus et relatifs) et proposer des itinéraires alternatifs robustes.
    
    \item \textbf{Résilience topologique et budgétaire du réseau de fibre optique :} 
    Comment planifier le déploiement ou l'extension de la fibre optique avec un budget d'investissement contraint, tout en garantissant des mécanismes de redondance et de raccordement de secours en cas de rupture physique d'un tronçon ? L'enjeu est de concrétiser un équilibre entre minimisation des coûts de construction et résistance structurelle face aux pannes.
\end{enumerate}

\subsection{Une démarche scientifique fondée sur l'aide à la décision}
Contrairement à une simple approche cartographique ou visuelle – souvent subjective et limitée aux représentations graphiques de surface –, la démarche adoptée dans ce projet s'inscrit dans une perspective scientifique rigoureuse. Chaque décision d'aménagement ou préconisation opérationnelle s'appuie sur :
\begin{itemize}
    \item Des \textbf{indicateurs quantitatifs formels} (taux d'accessibilité, degrés de connexité, coûts d'infrastructure, métriques de retard).
    \item Des \textbf{comparaisons algorithmiques approfondies} (évaluation croisée de la complexité et de l'efficacité des méthodes de parcours, de plus courts chemins et d'arbres couvrants).
    \item Des \textbf{analyses de robustesse} fondées sur la simulation de scénarios de perturbation réalistes appliqués sur des environnements isolés, garantissant la reproductibilité et la traçabilité des résultats d'ingénierie.
\end{itemize}